\newpage
\scp{\tsc{East Seram}}{11-02}
The \tsc{East Seram} node contains certain languages of eastern
Seram, minimally \ili{Bobot} and \ili{Masiwang}.
\ili{Salas} may also belong in this node,
but no available data is available to be sure (\srf{07-05}).
\citet{LoskiLoski1989} link it
with \tsc{\ili{Ambon}-Seram} languages to its west.

Following \citet{Collins1986}, the main change which unites \ili{Bobot}
and \ili{Masiwang} is medial *-ŋ- > \it{k}.
In addition, these two languages share *d/*z/*l/*j > \it{l};
*ma-p > \it{b}; and *k > {\0},
\it{ʔ} (preserved as \it{k-} in some verbs
and body parts that take prefixes).
\citet{Collins1986} also proposes *R > **ɣ as an exclusively shared innovation.
However, the modern reflexes (\ili{Bobot} *R > \it{w}, {\0}; \ili{Masiwang}
*R > \it{s}) mean that shared *R > **ɣ is somewhat speculative.

Of these changes, medial *-ŋ- > \it{k},
as well as the different reflexes of *R compared with *d/*z/*l/*j
clearly set these languages apart from \tsc{\ili{Ambon}-Seram}.
\citet{Collins1986} also includes \tsc{Seti} within \tsc{East Seram},
but we place it within \tsc{\ili{Ambon}-Seram} (\srf{10-03-04}) on
the basis of  *ŋ > \it{n},\footnote{
		There is one possible instance
		of *ŋ > \it{k} word medially in \tsc{Seti}; *naŋuy > \it{nak} ‘swim’,
		against three instances of medial *ŋ > \it{n}.
		Given *ŋ > \it{n} is the majority reflex, the word for ‘swim’
		in \tsc{Seti} is attributable to contact with neighbouring \ili{Bobot}.
		See \srf{10-03-04} for more discussion of \tsc{Seti}.}
as well as the usual merger of *d/*z/*l/*j > \it{l}.
\tsc{Seti} does have some instances of *R > \it{h} which may
be due to contact with \tsc{East Seram},
particularly if Collins’ proposal of *R > **ɣ is accepted.
\citeauthor{Collins1986}’ inclusion of \tsc{Seti} within \tsc{East Seram}
leads him to propose a subsequent node,
“East Rivers”, for \ili{Bobot} and \ili{Masiwang}.
However, under our classification this node is not necessary
as \ili{Bobot} and \ili{Masiwang} are the only languages
we can confidently place within \tsc{East Seram}.

Reflexes of *ŋ are given in \trf{11-08},
alongside cognates in other STB languages for comparison.
Word initially *ŋ- > \it{n}
and word medially *-ŋ- > \it{k}.
\ili{Masiwang} *naŋuy > \it{ke-nani} ‘swim’ has irregular medial *-ŋ- > \it{n}.
Medial *-ŋ- > \it{k} has parallels in several other languages/branches
of STB, though in these languages it usually also occurs word initially.
\ili{Kowiai} (\tsc{Eastern Islands}) has *ŋ > \it{g, ŋg} (\srf{11-03-02}),
\ili{Teor} has *ŋ > \it{g} (\srf{11-04})
\tsc{Bomberai Tip} has *ŋ > \it{g, ŋg} (\srf{11-05-01}), and
\tsc{East Bomberai} has *ŋ > **g (\srf{11-05-02}).

\begin{table}\tcp{\tsc{East Seram} *ŋ > \it{n}-, -\it{k}-}{11-08}
\begin{adjustbox}{width=\textwidth}
	\begin{threeparttable}\st{2pt}
	\begin{tabular}{llllllllll}\lsptoprule
gloss	&	name	&	year\su{†}	&	swim	&	wind	&	hear	&	ear	&	sky	&	cry	&	branch\\
PMP	&	*\tbr{ŋ}ajan	&	**\tbr{ŋ}aRək	&	*na\tbr{ŋ}uy	&	*ha\tbr{ŋ}in	&	*də\tbr{ŋ}əR	&	*tali\tbr{ŋ}a	&	*la\tbr{ŋ}it	&	*ta\tbr{ŋ}is	&	*sa\tbr{ŋ}a\\	\midrule
\ili{Bobot}	&	-\tbr{n}alan	&	{\hr}\tbr{n}aa	&	{\hr}na\tbr{k}u̥	&	{\hr}ya\tbr{k}in	&	{\hr}blo\tbr{k}a	&	{\hr}tali\tbr{k}a-n	&	{\hr}la\tbr{k}it	&	{\hr}ta\tbr{k}it	&	{\hr}ai ya\tbr{k}an\\
\ili{Masiwang}	&	{\hr}\tbr{n}ala-n	&	{\hr}\tbr{n}alaʔ	&	{\hr}ke-na\tbr{n}i	&	{\hr}ta\tbr{k}in	&	{\hr}lo\tbr{k}a	&	{\hr}tali\tbr{k}a-n	&	{\hr}la\tbr{k}it	&	{\hr}ra\tbr{k}it	&	{\hr}sak{\tl}sa\tbr{k}an\\
\ili{Geser}	&	{\hr}ŋasan	&	{\hr}ŋarak	&	{\hr}naŋu	&	{\hr}aŋin	&	{\hr}roŋan	&	{\hr}tiliŋa	&	{\hr}laŋit	&	\cg{(sela)}	&	{\hr}saŋ\\
\ili{Yamdena}	&	{\hr}ŋara-ny	&	{\hr}ŋarak	&	\cg{(n-luri)}	&	\cg{(mnaur)}	&	{\hr}-deŋar	&	{\hr}tliŋa-ny	&	{\hr}laŋit	&	{\hr}-tasiŋ	&	{\hr}saŋa-ny\\
\ili{Kei}	&	\cg{(mema-)}	&	\cg{(taun)}	&	{\hr}snaŋ	&	\cg{(niet)}	&	{\hr}denar	&	\cg{(aru-n)}	&	{\hr}lan	&	{\hr}tnanit	&	{\hr}haŋa\\
		\lspbottomrule
	\end{tabular}
		\begin{tablenotes}\tnsize
			\item[†] {Compare \ili{Buru} \it{ŋahe-t},
								\ili{Asilulu} \it{nale}, East \ili{Kola} \it{nahak},
								Southwest \ili{Tarangan} \it{narak}
								all ‘year’.}
		\end{tablenotes}
	\end{threeparttable}
\end{adjustbox}
\end{table}

%\begin{table}\tcp{\tsc{East Seram} *ŋ > \it{n}-, -\it{k}-}{11-08}
%%\begin{adjustbox}{width=\textwidth}
	%\begin{threeparttable}{\st{6pt}
	%\begin{tabular}{llllll}\lsptoprule
%gloss	&	name	&	year\su{†}	&	swim	&	wind	&	hear	\\
%PMP	&	*\tbr{ŋ}ajan	&	**\tbr{ŋ}aRək	&	*na\tbr{ŋ}uy	&	*ha\tbr{ŋ}in	&	*də\tbr{ŋ}əR	\\	\midrule
%\ili{Bobot}	&	-\tbr{n}alan	&	{\hr\hr}\tbr{n}aa	&	{\hr}na\tbr{k}u̥	&	{\hr}ya\tbr{k}in	&	{\hr}a-vlo\tbr{k}a	\\
%\ili{Masiwang}	&	{\hr}\tbr{n}ala-n	&	{\hr\hr}\tbr{n}alaʔ	&	{\hr}ke-na\tbr{n}i	&	{\hr}ta\tbr{k}in	&	{\hr}lo\tbr{k}a	\\	\midrule
%\ili{Geser}	&	{\hr}ŋasan	&	{\hr\hr}ŋarak	&	{\hr}naŋu	&	{\hr}aŋin	&	{\hr}roŋan	\\
%\ili{Yamdena}	&	{\hr}ŋara-ny	&	{\hr\hr}ŋarak	&	\cg{(n-luri)}	&	\cg{(mnaur)}	&	{\hr}-deŋar	\\
%\ili{Kei}	&	\cg{(mema-)}	&	\hr\cg{(taun)}	&	{\hr}snaŋ	&	\cg{(niet)}	&	{\hr}denar	\\
		%\midrule\midrule
	%\end{tabular}}
	%{\st{6pt}
	%\begin{tabular}{lllll}
%gloss	&	ear	&	sky	&	cry	&	branch\\
%PMP	&	*tali\tbr{ŋ}a	&	*la\tbr{ŋ}it	&	*ta\tbr{ŋ}is	&	*sa\tbr{ŋ}a\\	\midrule
%\ili{Bobot}	&	{\hr}tali\tbr{k}a-n	&	{\hr}la\tbr{k}it	&	{\hr}ta\tbr{k}it	&	{\hr}ai ya\tbr{k}an\\
%\ili{Masiwang}	&	{\hr}tali\tbr{k}a-n	&	{\hr}la\tbr{k}it	&	{\hr}ra\tbr{k}it	&	{\hr}sak{\tl}sa\tbr{k}an\\	\midrule
%\ili{Geser}	&	{\hr}tiliŋa	&	{\hr}laŋit	&	\cg{(sela)}	&	{\hr}saŋ\\
%\ili{Yamdena}	&	{\hr}tliŋa-ny	&	{\hr}laŋit	&	{\hr}-tasiŋ	&	{\hr}saŋa-ny\\
%\ili{Kei}	&	\cg{(aru-n)}	&	{\hr}lan	&	{\hr}tnanit	&	{\hr}haŋa\\
		%\lspbottomrule
	%\end{tabular}}
		%\begin{tablenotes}\tnsize
			%\item[†] {Compare \ili{Buru} \it{ŋahet},
								%\ili{Asilulu} \it{nale}, East \ili{Kola} \it{nahak},
								%Southwest \ili{Tarangan} \it{narak}
								%all ‘year’.}
		%\end{tablenotes}
	%\end{threeparttable}
%%\end{adjustbox}
%\end{table}

The reflexes of *R in \tsc{East Seram} are unique among languages of the region.
\ili{Bobot} has *R > {\0} \it{/}V\tsc{[+back]}
and *R > \it{w} elsewhere.
\ili{Masiwang} has *R > \it{s},
thus merging with *s.
Examples are given in \trf{11-09},
alongside reflexes in other STB languages for comparison.

\begin{table}\tcp{\tsc{East Seram} *R}{11-09}
\begin{adjustbox}{width=\textwidth}
	\st{2pt}
	\begin{tabular}{llllllll}\lsptoprule
local gl.	&	house	&	blood	&	rope	&	bone	&	new	&	man	&	bite\\
PMP	&	*\tbr{R}umaq	&	*da\tbr{R}aq	&	*wa\tbr{R}əj	&	*du\tbr{R}i	&	*baqə\tbr{R}u	&	*ma\tbr{R}uqanay	&	*ka\tbr{R}at\\	\midrule
\ili{Bobot}	&	{\hr}uma	&	{\hr}la\tbr{w}a	&	{\hr}wa\tbr{w}at	&	{\hr}-luin	&	{\hr}vo{\tl}voi	&		&	\\
\ili{Masiwang}	&	{\hr}\tbr{s}uma	&	{\hr}la\tbr{s}a	&	{\hr}wa\tbr{s}at	&	{\hr}lu\tbr{s}i-n	&	{\hr}hosi{\tl}ho\tbr{s}in	&	{\hr}me\tbr{s}anem	&	{\hr}ka\tbr{s}at\\	\midrule
\ili{Geser}	&	{\hr}ruma	&	{\hr}rara	&	\cg{(tali)}	&	{\hr}ruri	&	{\hr}wou{\tl}wou	&	{\hr}urana	&	{\hr}karat\\
\ili{Yamdena}	&	\cg{(rahan)}	&	{\hr}dara-ny	&	{\hr}waras	&	{\hr}duri	&	{\hr}be{\tl}bery	&	{\hr}merwan-e	&	{\hr}n-karat\\
\ili{Kei}	&	\cg{(raha)}	&	{\hr}lar	&	{\hr}warat	&	{\hr}luri	&	{\hr}virun	&	{\hr}bran	&	\cg{(kik)}\\
		\lspbottomrule
	\end{tabular}
\end{adjustbox}
\end{table}

The reflexes of *R in \tsc{East Seram} are distinct from those
of *d/*z/*l/*j, the latter of which merge to \it{l}.
This is different from all \tsc{\ili{Ambon}-Seram} languages (except
Boano) in which *l/*R/*j > **l.
On the other hand,
the merger of *l/*j > \it{l} sets \tsc{East Seram} apart from
other STB languages as *l
and *j are kept distinct in most other branches in most words.
The merger of *l/*j also occurs in \ili{Banda} (\srf{12-02}).
Examples of *l/*j > \it{l} are given in \trf{11-10}.
And examples of *l = \it{l} in \trf{11-10b}.
Examples of *d > \it{l}
and *z > \it{l} were given in \trf{11-02} and \trf{11-03} respectively.

Apart from the usual *l = \it{l,} there are also a handful of
irregular reflexes of *l in the available data: *\tbr{l}ima > \ili{Bobot}
\it{-\tbr{n}iman} ‘arm’, *qu\tbr{l}əj > \ili{Bobot} \it{u\tbr{r}at} ‘worm’,
*kali > \ili{Masiwang} \it{kai\tbr{s}} ‘dig’ (with final CV > VC metathesis),
and *ku\tbr{l}it > \ili{Bobot} \it{-ku\tbr{r}it}, \ili{Masiwang} \it{u\tbr{h}it} ‘skin’.
There is also one possible case of *j > \it{n}; \ili{Masiwang} *ŋi\tbr{j}uŋ
> \it{ni\tbr{n}in}, as well as one possible case of *j > {\0} in *qalejaw > \ili{Bobot}
\it{lia matan} ‘sun’ (\ili{Masiwang} \it{fulan} ‘sun’),
though the regularity of *j > {\0} in this word in \tsc{\ili{Ambon}-Seram}
(\srf{07-03-02}) indicates that this may be a loan.

\begin{table}
	\centering\tcp{\tsc{East Seram} *j > \it{l}}{11-10}
	\st{6pt}\begin{tabular}{lllllll}\lsptoprule
*gloss	&	how m.	&	name	&	dry	&	nose	&	ySib	&	gall\\
PMP	&	*pi\tbr{j}a	&	*ŋa\tbr{j}an	&	*ma\tbr{j}a	&	*\cg{(ŋ)}i\tbr{j}uŋ	&	*hua\tbr{j}i	&	*qapə\tbr{j}u\\	\midrule
\ili{Bobot}	&	{\hr}fi\tbr{l}a	&	-na\tbr{l}an	&	{\hr}ma-ma\tbr{l}a	&	{\hr}-i\tbr{l}in	&	{\hr}wa\tbr{l}i-n	&	{\hr}fo\tbr{l}i\\
\ili{Masiwang}	&	{\hr}hi\tbr{l}a	&	{\hr}na\tbr{l}a-n	&	{\hr}mo-ma\tbr{l}a	&	{\hr}ni\tbr{n}in	&	{\hr}wa\tbr{l}i(-n)	&	{\hr}ho\tbr{l}i\\	\midrule
\ili{Geser}	&	{\hr}fis	&	{\hr}ŋasan	&	{\hr}kirimas	&	{\hr}so	&	{\hr}ali	&	{\hr}foli\\
\ili{Yamdena}	&	{\hr}fir	&	{\hr}ŋara-ny	&	\cg{(-ŋretw)}	&	{\hr}iri-ŋw	&	{\hr}wai-ny	&	{\hr}feru\\
\ili{Onin}	&	{\hr}firas	&	{\hr}gara-r	&	\cg{(atetar)}	&	{\hr}iri-r	&	{\hr}wari-r	&	{\hr}marik\\
		\lspbottomrule
	\end{tabular}
\end{table}

\begin{table}
	\centering\tcp{\tsc{East Seram} *l = \it{l}}{11-10b}
	\st{6pt}\begin{tabular}{llllll}\lsptoprule
*gloss	&	moon	&	laugh	&	three	&	head	&	egg	\\
PMP	&	*bu\tbr{l}an	&	**ma\tbr{l}ip	&	*tə\tbr{l}u	&	*qu\tbr{l}u	&	*qatə\tbr{l}uR	\\ \midrule
\ili{Bobot}	&	{\hr}vu\tbr{l}an	&	{\hr}ba\tbr{l}íʔ	&	{\hr}to\tbr{l}(u̥)	&	-u\tbr{l}in katin	&	{\hr}to\tbr{l}in	\\
\ili{Masiwang}	&	{\hr}hu\tbr{l}an	&	{\hr}m\tbr{l}i	&	{\hr}ain-to\tbr{l}i	&	{\hr}u\tbr{l}in	&	{\hr}to\tbr{l}in	\\	\midrule
\ili{Geser}	&	{\hr}ulan	&	{\hr}malif	&	{\hr}tolu	&	{\hr}lu	&	{\hr}tolu	\\
\ili{Yamdena}	&	{\hr}bulan	&	{\hr}n-malip	&	{\hr}tel	&	{\hr}ulu	&	\cg{(batin)}	\\
\ili{Onin}	&	{\hr}punan	&	{\hr}a-manif	&	{\hr}teni	&	{\hr}unin	&	\cg{(patin)}	\\
		\lspbottomrule
	\end{tabular}
\end{table}

The vowel *u undergoes changes in final syllables in \tsc{East Seram}.
Before a final consonant both languages usually have *u/*i > \it{i}.
In final open syllables \ili{Masiwang} usually has *u/*i > \it{i,}
while \ili{Bobot} has reduction of *u/*i.
In the \ili{Atiahu} dialect of \ili{Bobot} the usual change is *u > {\0}
in all cases except after /h/ where it is retained as \it{u} \citep[138]{Collins1986}.
In the Werinama dialect,
from which most data in \citet{Collins1986} comes,
*u/*i > {\0} in most cases after a voiced consonant,
while after a voiceless consonant *u/*i is usually reduced to a voiceless vowel.
This is transcribed by \citet{Collins1986} as ‹u̥› in the main
text and as ‹u̥› or ‹V̥› in the accompanying word list.
This ‹V̥› is described as:
“[{\ldots}] an unspecified devoiced vowel” \citep[141]{Collins1986}.

\citet{Collins1986} also observes that final *i metathesises
with word-final *n in \ili{Masiwang}.
We have identified one possible additional example *kali > \it{kais}
‘dig’ (though this word would also have irregular *l > \it{s}).
Metathesis of high vowels with a previous consonant is an areal
feature of eastern Seram
and also occurs in \tsc{Seti} and \ili{Bati}.
Similarly, the change of final *u > \it{i} occurs in nearby
\tsc{Bomberai Tip} as well as \ili{Yamdena}.\footnote{
		Final *u > \it{i} also occurs
		in \ili{Wetan} and \tsc{Babar} (\tsc{Southwest Maluku}, \tsc{Timor-Babar}; \srf{04-06}).}
Examples of \tsc{East Seram} syllable-final
*u and *i before a final consonant are given in \trf{11-11}
and examples in final open syllables are given in \trf{11-12}.
“\ili{Bobot} (We.)” is from the Werinama dialect of \ili{Bobot}
and “\ili{Bobot} (At.)” is from the \ili{Atiahu} dialect.
\ili{Atiahu} \ili{Bobot} data is from \citet{Wallace1869},
with the exception of \it{hat} ‘stone’ from \citet[141]{Collins1986}.

\begin{table}
	\centering\tcp{\tsc{East Seram} *u and *i /{\gap}C{\#}}{11-11}
\begin{adjustbox}{width=\textwidth}
	\st{2pt}\begin{tabular}{lllllllll}\lsptoprule
local gl.	&	br\ili{east}	&	head	&	fear	&	egg	&	louse	&	feather	&	mouth	&	sea\\
PMP	&	*sus\tbr{u}	&	*qul\tbr{u}	&	*tak\tbr{u}t	&	*qatəl\tbr{u}R	&	*kut\tbr{u}	&	*bul\tbr{u}	&	*bib\tbr{i}R	&	*tas\tbr{i}k\\	\midrule
\ili{Bobot} (We.)	&	-sus\tbr{i}n	&	-ul\tbr{i}n katin	&	{\hr}mta\tbr{i}t	&	{\hr}tol\tbr{i}n	&	{\hr}kut\tbr{i}n	&	{\hr}vul\tbr{i}n	&	{\hr}viv\tbr{i}n	&	{\hr}tas\tbr{i}\\
\ili{Bobot} (At.)	&		&	‹yúl\tbr{i}n›	&		&	‹tól\tbr{i}n›	&	‹ut\tbr{u}›	&	‹ful\tbr{i}n›	&	\cg{(‹vudin›)}	&	‹tás\tbr{i}›\\
\ili{Masiwang}	&	{\hr}sus\tbr{i}n	&	{\hr}ul\tbr{i}n	&	{\hr}kato\tbr{i}t	&	{\hr}tol\tbr{i}n	&	{\hr}taniman	&	{\hr}hul{\tl}hul\tbr{i}n	&	\cg{(uban)}	&	{\hr}tas\tbr{i}ʔ\\
		\lspbottomrule
	\end{tabular}
\end{adjustbox}
\end{table}

\begin{table}
	\centering\tcp{\tsc{East Seram} *u and *i /{\gap}{\#}}{11-12}
\begin{adjustbox}{width=\textwidth}
	\st{2pt}\begin{tabular}{llllllllll}\lsptoprule
local gl.	&	stone	&	dog	&	pound	&	burn	&	kill	&	fire	&	drink	&	white	&	dig\\
PMP	&	*bat\tbr{u}	&	*as\tbr{u}	&	*tukt\tbr{u}k	&	*tun\tbr{u}	&	*bun\tbr{u}q	&	*hap\tbr{u}y	&	*in\tbr{u}m	&	*ma-put\tbr{i}q	&	*kal\tbr{i}\\	\midrule
\ili{Bobot} (We.)	&	{\hr}vatu̥	&	{\hr}yasu̥	&	{\hr}dutu̥	&	{\hr}dun	&	{\hr}vun(V̥)	&	{\hr}yafV̥	&	{\hr}ninV̥	&	{\hr}babutu̥	&	{\hr}kalV̥\\
\ili{Bobot} (At.)	&	{\hr}hat	&	‹yás›	&		&		&		&	‹yaf›	&		&	‹babút›	&	\\
\ili{Masiwang}	&	{\hr}hat\tbr{i}	&	{\hr}ya\tbr{i}	&	{\hr}tut\tbr{i}	&	{\hr}tu\tbr{i}n	&	{\hr}hu\tbr{i}n	&	{\hr}yah	&	{\hr}nim\tbr{i}n	&	{\hr}but\tbr{i}	&	{\hr}ka\tbr{i}s\\
		\lspbottomrule
	\end{tabular}
\end{adjustbox}
\end{table}
